% GENERATED FILE. Edit canonical lessons, then run npm run generate:books.
% canonical-source-hash: fnv1a64:c49ff743bdaa9a35
% canonical-lessons: TA-C294-catci, TA-C294-campavam, TA-C294-tirutu, TA-C294-kollai, TA-C294-nikal

\chapter{A Witness and a Theft}
\label{ch:ta-a-witness-and-a-theft}
\hlchaptermodality{294}

\begin{chapteropening}
\textbf{By the end of this chapter:} \emph{I can talk about a theft and the people who saw it.}

Complete the last lesson of chapter 294: I can talk about a theft and the people who saw it.
\end{chapteropening}

\section[\texorpdfstring{\ta{சாட்சி} (cāṭci)}{சாட்சி (cāṭci)}]{\ta{சாட்சி} (cāṭci) — a witness}

\label{lesson:TA-C294-catci}



\begin{quote}
Before the new one: say the Tamil for a complaint, then the Tamil for an athlete.
\end{quote}

\subsection*{You'll want to know: \ta{சாட்சி}}
\textbf{\ta{சாட்சி}} — \emph{cāṭci} — \textquotedblleft{}a witness\textquotedblright{}.

\begin{quote}
\ta{எல்லாம்} \ta{சாட்சியைப்} \ta{பொறுத்தது}.
\end{quote}

\begin{quote}
\emph{ellām cāṭciyaip poṟuttatu.}
\end{quote}

\begin{quote}
Everything depends on the witness.
\end{quote}

Before -\ta{ஐ} the final -\ta{இ} takes the glide \ta{ய்}: \ta{சாட்சியை}.

\begin{cousinweb}
\textbf{\ta{சாட்சி}} (\emph{cāṭci}) is borrowed from Sanskrit \emph{sākṣin}, \textquotedblleft{}a witness\textquotedblright{} (CDIAL 13321); Tamil writes the Sanskrit kṣ as \ta{ட்ச}.
\end{cousinweb}

\begin{grammarlens}[title={What you've built}]
One more word for telling what happened.
\end{grammarlens}

\subsection*{Your turn}
\begin{itemize}
\raggedright
  \item \emph{Say it:} \emph{cāṭci}
  \item \emph{Say it:} \emph{cāṭci}, once more
  \item \emph{Say it:} \emph{ellām cāṭciyaip poṟuttatu.} — Everything depends on the witness.
  \item \emph{Recall:} say the Tamil for a complaint, then the Tamil for an athlete, then say \emph{cāṭci} again
\end{itemize}

\begin{tcolorbox}[breakable,colback=teal!4,colframe=teal!35!black,title={Before you move on}]
What does \ta{சாட்சி} mean? (\textquotedblleft{}A witness\textquotedblright{}.) Say it once more.
\end{tcolorbox}
\section[\texorpdfstring{\ta{சம்பவம்} (campavam)}{சம்பவம் (campavam)}]{\ta{சம்பவம்} (campavam) — an incident}

\label{lesson:TA-C294-campavam}



\begin{quote}
Before the new one: say the Tamil for an athlete, then the Tamil for a witness.
\end{quote}

\subsection*{You'll want to know: \ta{சம்பவம்}}
\textbf{\ta{சம்பவம்}} — \emph{campavam} — \textquotedblleft{}an incident\textquotedblright{}.

\begin{quote}
\ta{நேற்று} \ta{இரவு} \ta{எங்கள்} \ta{தெருவில்} \ta{ஒரு} \ta{சம்பவம்} \ta{நடந்தது}.
\end{quote}

\begin{quote}
\emph{nēṟṟu iravu eṅkaḷ teruvil oru campavam naṭantatu.}
\end{quote}

\begin{quote}
Last night an incident happened on our street.
\end{quote}

Plural: \ta{சம்பவங்கள்}, \emph{campavaṅkaḷ}.

A -\ta{அம்} noun takes -\ta{அத்து}- before an ending: \ta{சம்பவத்தை}, \ta{சம்பவத்தில்}.

\begin{cousinweb}
\textbf{\ta{சம்பவம்}} (\emph{campavam}) is borrowed from Sanskrit \emph{sambhava}, \textquotedblleft{}coming into being, an occurrence\textquotedblright{} (Monier-Williams); the Sanskrit final -a becomes Tamil -\ta{அம்}, and bh is written \ta{ப}.
\end{cousinweb}

\begin{grammarlens}[title={What you've built}]
One more word for telling what happened.
\end{grammarlens}

\subsection*{Your turn}
\begin{itemize}
\raggedright
  \item \emph{Say it:} \emph{campavam}
  \item \emph{Say it:} \emph{campavam}, once more
  \item \emph{Say it:} \emph{nēṟṟu iravu eṅkaḷ teruvil oru campavam naṭantatu.} — Last night an incident happened on our street.
  \item \emph{Recall:} say the Tamil for an athlete, then the Tamil for a witness, then say \emph{campavam} again
\end{itemize}

\begin{tcolorbox}[breakable,colback=teal!4,colframe=teal!35!black,title={Before you move on}]
What does \ta{சம்பவம்} mean? (\textquotedblleft{}An incident\textquotedblright{}.) Say it once more.
\end{tcolorbox}
\section[\texorpdfstring{\ta{திருடு} (tiruṭu)}{திருடு (tiruṭu)}]{\ta{திருடு} (tiruṭu) — to steal}

\label{lesson:TA-C294-tirutu}



\begin{quote}
Before the new one: say the Tamil for a witness, then the Tamil for an incident.
\end{quote}

\subsection*{You'll want to know: \ta{திருடு}}
\textbf{\ta{திருடு}} — \emph{tiruṭu} — \textquotedblleft{}to steal\textquotedblright{}.

\begin{quote}
\ta{அவன்} \ta{என்} \ta{பையைத்} \ta{திருடினான்}.
\end{quote}

\begin{quote}
\emph{avaṉ eṉ paiyait tiruṭiṉāṉ.}
\end{quote}

\begin{quote}
He stole my bag.
\end{quote}

Past (I …): \ta{திருடினேன்}, \emph{tiruṭiṉēṉ}.

What is stolen takes -\ta{ஐ}: \ta{பையைத்} \ta{திருடு}, to steal a bag.

\begin{cousinweb}
\textbf{\ta{திருடு}} (\emph{tiruṭu}) is native Dravidian (DEDR 3250); the same entry gives \ta{திருடன்} (\emph{tiruṭaṉ}), \textquotedblleft{}a thief\textquotedblright{}.
\end{cousinweb}

\begin{grammarlens}[title={What you've built}]
One more word for telling what happened.
\end{grammarlens}

\subsection*{Your turn}
\begin{itemize}
\raggedright
  \item \emph{Say it:} \emph{tiruṭu}
  \item \emph{Say it:} \emph{tiruṭu}, once more
  \item \emph{Say it:} \emph{avaṉ eṉ paiyait tiruṭiṉāṉ.} — He stole my bag.
  \item \emph{Recall:} say the Tamil for a witness, then the Tamil for an incident, then say \emph{tiruṭu} again
\end{itemize}

\begin{tcolorbox}[breakable,colback=teal!4,colframe=teal!35!black,title={Before you move on}]
What does \ta{திருடு} mean? (\textquotedblleft{}To steal\textquotedblright{}.) Say it once more.
\end{tcolorbox}
\section[\texorpdfstring{\ta{கொள்ளை} (koḷḷai)}{கொள்ளை (koḷḷai)}]{\ta{கொள்ளை} (koḷḷai) — a robbery; plunder}

\label{lesson:TA-C294-kollai}



\begin{quote}
Before the new one: say the Tamil for an incident, then the Tamil for to steal.
\end{quote}

\subsection*{You'll want to know: \ta{கொள்ளை}}
\textbf{\ta{கொள்ளை}} — \emph{koḷḷai} — \textquotedblleft{}a robbery; plunder\textquotedblright{}.

\begin{quote}
\ta{அந்த} \ta{வங்கியில்} \ta{பெரிய} \ta{கொள்ளை} \ta{நடந்தது}.
\end{quote}

\begin{quote}
\emph{anta vaṅkiyil periya koḷḷai naṭantatu.}
\end{quote}

\begin{quote}
A big robbery took place at that bank.
\end{quote}

A theft by force or on a large scale; \ta{கொள்ளை} \ta{நடந்தது}, a robbery took place.

\begin{cousinweb}
\textbf{\ta{கொள்ளை}} (\emph{koḷḷai}) is native Dravidian: DEDR 2151 lists it beside \ta{கொள்}, \textquotedblleft{}to seize, take\textquotedblright{}, so a robbery is at root a seizing.
\end{cousinweb}

\begin{grammarlens}[title={What you've built}]
One more word for telling what happened.
\end{grammarlens}

\subsection*{Your turn}
\begin{itemize}
\raggedright
  \item \emph{Say it:} \emph{koḷḷai}
  \item \emph{Say it:} \emph{koḷḷai}, once more
  \item \emph{Say it:} \emph{anta vaṅkiyil periya koḷḷai naṭantatu.} — A big robbery took place at that bank.
  \item \emph{Recall:} say the Tamil for an incident, then the Tamil for to steal, then say \emph{koḷḷai} again
\end{itemize}

\begin{tcolorbox}[breakable,colback=teal!4,colframe=teal!35!black,title={Before you move on}]
What does \ta{கொள்ளை} mean? (\textquotedblleft{}A robbery; plunder\textquotedblright{}.) Say it once more.
\end{tcolorbox}
\section[\texorpdfstring{\ta{நிகழ்} (nikaḻ)}{நிகழ் (nikaḻ)}]{\ta{நிகழ்} (nikaḻ) — to happen; to occur}

\label{lesson:TA-C294-nikal}



\begin{quote}
Before the new one: say the Tamil for to steal, then the Tamil for a robbery.
\end{quote}

\subsection*{You'll want to know: \ta{நிகழ்}}
\textbf{\ta{நிகழ்}} — \emph{nikaḻ} — \textquotedblleft{}to happen; to occur\textquotedblright{}.

\begin{quote}
\ta{இந்தச்} \ta{சம்பவம்} \ta{நேற்று} \ta{மாலை} \ta{நிகழ்ந்தது}.
\end{quote}

\begin{quote}
\emph{intac campavam nēṟṟu mālai nikaḻntatu.}
\end{quote}

\begin{quote}
This incident happened yesterday evening.
\end{quote}

Past (it …): \ta{நிகழ்ந்தது}, \emph{nikaḻntatu}.

A written word: in speech \ta{நடந்தது} says the same.

\begin{cousinweb}
\textbf{\ta{நிகழ்}} (\emph{nikaḻ}) is native Dravidian (DEDR 3661); the same entry gives \ta{நிகழ்} \ta{காலம்}, the present tense, \textquotedblleft{}the time that is happening\textquotedblright{}.
\end{cousinweb}

\begin{grammarlens}[title={What you've built}]
One more word for telling what happened.
\end{grammarlens}

\subsection*{Your turn}
\begin{itemize}
\raggedright
  \item \emph{Say it:} \emph{nikaḻ}
  \item \emph{Say it:} \emph{nikaḻ}, once more
  \item \emph{Say it:} \emph{intac campavam nēṟṟu mālai nikaḻntatu.} — This incident happened yesterday evening.
  \item \emph{Recall:} say the Tamil for to steal, then the Tamil for a robbery, then say \emph{nikaḻ} again
\end{itemize}

\begin{tcolorbox}[breakable,colback=teal!4,colframe=teal!35!black,title={Before you move on}]
What does \ta{நிகழ்} mean? (\textquotedblleft{}To happen; to occur\textquotedblright{}.) Say it once more.
\end{tcolorbox}
